%=============================================================================
% Chapter 3 Exercises
%=============================================================================
\newpage
\section{Chapter 3 Exercises [725]}

%-----------------------------------------------------------------------------
\subsection{[5] Unsigned Hex Subtraction}
\hypertarget{prob-3.1}{}
\noindent$\langle$\S3.2$\rangle$\medskip

\noindent What is $\texttt{5ED4} - \texttt{07A4}$ when these values represent
unsigned 16-bit hexadecimal numbers? The result should be written in
hexadecimal. Show your work.

\problembreak

%-----------------------------------------------------------------------------
\subsection{[5] Signed Hex Subtraction (Sign-Magnitude)}
\hypertarget{prob-3.2}{}
\noindent$\langle$\S3.2$\rangle$\medskip

\noindent What is $\texttt{5ED4} - \texttt{07A4}$ when these values represent
signed 16-bit hexadecimal numbers stored in sign-magnitude format? The result
should be written in hexadecimal. Show your work.

\problembreak

%-----------------------------------------------------------------------------
\subsection{[10] Hex to Binary Conversion}
\hypertarget{prob-3.3}{}
\noindent$\langle$\S3.2$\rangle$\medskip

\noindent Convert \texttt{5ED4} into a binary number. What makes base 16
(hexadecimal) an attractive numbering system for representing values in
computers?

\problembreak

%-----------------------------------------------------------------------------
\subsection{[5] Unsigned Octal Subtraction}
\hypertarget{prob-3.4}{}
\noindent$\langle$\S3.2$\rangle$\medskip

\noindent What is $4365 - 3412$ when these values represent unsigned 16-bit
octal numbers? The result should be written in octal. Show your work.

\problembreak

%-----------------------------------------------------------------------------
\subsection{[5] Signed Octal Subtraction (Sign-Magnitude)}
\hypertarget{prob-3.5}{}
\noindent$\langle$\S3.2$\rangle$\medskip

\noindent What is $4365 - 3412$ when these values represent signed 12-bit
numbers stored in sign-magnitude format? The result should be written in
octal. Show your work.

\problembreak

%-----------------------------------------------------------------------------
\subsection{[5] Unsigned 8-bit Subtraction Overflow}
\hypertarget{prob-3.6}{}
\noindent$\langle$\S3.2$\rangle$\medskip

\noindent Assume 185 and 122 are unsigned 8-bit decimal integers. Calculate
$185 - 122$. Is there overflow, underflow, or neither?

\problembreak

%-----------------------------------------------------------------------------
\subsection{[5] Signed Sign-Magnitude Subtraction Overflow}
\hypertarget{prob-3.7}{}
\noindent$\langle$\S3.2$\rangle$\medskip

\noindent Assume 185 and 122 are signed 8-bit decimal integers stored in
sign-magnitude format. Calculate $185 - 122$. Is there overflow, underflow,
or neither?

\problembreak

%-----------------------------------------------------------------------------
\subsection{[5] Signed Sign-Magnitude Addition Overflow}
\hypertarget{prob-3.8}{}
\noindent$\langle$\S3.2$\rangle$\medskip

\noindent Assume 185 and 122 are signed 8-bit decimal integers stored in
sign-magnitude format. Calculate $185 + 122$. Is there overflow, underflow,
or neither?

\problembreak

%-----------------------------------------------------------------------------
\subsection{[10] Two's Complement Saturating Addition}
\hypertarget{prob-3.9}{}
\noindent$\langle$\S3.2$\rangle$\medskip

\noindent Assume 151 and 214 are signed 8-bit decimal integers stored in
two's complement format. Calculate $151 + 214$ using saturating arithmetic.
The result should be written in decimal. Show your work.

\problembreak

%-----------------------------------------------------------------------------
\subsection{[10] Two's Complement Saturating Subtraction}
\hypertarget{prob-3.10}{}
\noindent$\langle$\S3.2$\rangle$\medskip

\noindent Assume 151 and 214 are signed 8-bit decimal integers stored in
two's complement format. Calculate $151 - 214$ using saturating arithmetic.
The result should be written in decimal. Show your work.

\problembreak

%-----------------------------------------------------------------------------
\subsection{[10] Unsigned Saturating Addition}
\hypertarget{prob-3.11}{}
\noindent$\langle$\S3.2$\rangle$\medskip

\noindent Assume 151 and 214 are unsigned 8-bit integers. Calculate
$151 + 214$ using saturating arithmetic. The result should be written in
decimal. Show your work.

\problembreak

%-----------------------------------------------------------------------------
\subsection{[20] Multiplication Table (6-bit, Fig.~3.3)}
\hypertarget{prob-3.12}{}
\noindent$\langle$\S3.3$\rangle$\medskip

\noindent Using a table similar to that shown in Figure 3.6, calculate the
product of the octal unsigned 6-bit integers 62 and 12 using the hardware
described in Figure 3.3. You should show the contents of each register on
each step.

\problembreak

%-----------------------------------------------------------------------------
\subsection{[20] Multiplication Table (8-bit, Fig.~3.5)}
\hypertarget{prob-3.13}{}
\noindent$\langle$\S3.3$\rangle$\medskip

\noindent Using a table similar to that shown in Figure 3.6, calculate the
product of the octal unsigned 8-bit integers 62 and 12 using the hardware
described in Figure 3.5. You should show the contents of each register on
each step.

\problembreak

%-----------------------------------------------------------------------------
\subsection{[10] Multiply Time Calculation}
\hypertarget{prob-3.14}{}
\noindent$\langle$\S3.3$\rangle$\medskip

\noindent Calculate the time necessary to perform a multiply using the
approach given in Figures 3.3 and 3.5 if an integer is 8 bits wide and each
step of the operation takes 4 time units. Assume that in step 1a an addition
is always performed---either the multiplicand will be added, or a zero will
be. Also assume that the registers have already been initialized (you are
just counting how long it takes to do the multiplication loop itself). If
this is being done in hardware, the shifts of the multiplicand and multiplier
can be done simultaneously. If this is being done in software, they will have
to be done one after the other. Solve for each case.

\problembreak

%-----------------------------------------------------------------------------
\subsection{[10] Multiply Time with Stacked Adders}
\hypertarget{prob-3.15}{}
\noindent$\langle$\S3.3$\rangle$\medskip

\noindent Calculate the time necessary to perform a multiply using the
approach described in the text (31 adders stacked vertically) if an integer is
8 bits wide and an adder takes 4 time units.

\problembreak

%-----------------------------------------------------------------------------
\subsection{[20] Multiply Time (Fig.~3.7)}
\hypertarget{prob-3.16}{}
\noindent$\langle$\S3.3$\rangle$\medskip

\noindent Calculate the time necessary to perform a multiply using the
approach given in Figure 3.7 if an integer is 8 bits wide and an adder takes
4 time units.

\problembreak

%-----------------------------------------------------------------------------
\subsection{[20] Shift-and-Add Multiplication}
\hypertarget{prob-3.17}{}
\noindent$\langle$\S3.3$\rangle$\medskip

\noindent As discussed in the text, one possible performance enhancement is to
do a shift and add instead of an actual multiplication. Since $9 \times 6$,
for example, can be written $(2 \times 2 \times 2 + 1) \times 6$, we can
calculate $9 \times 6$ by shifting 6 to the left 3 times and then adding 6
to that result. Show the best way to calculate
$\texttt{0x33} \times \texttt{0x55}$ using shifts and adds/subtracts. Assume
both inputs are unsigned 8-bit integers.

\problembreak

%-----------------------------------------------------------------------------
\subsection{[20] Division Table (Fig.~3.8)}
\hypertarget{prob-3.18}{}
\noindent$\langle$\S3.4$\rangle$\medskip

\noindent Using a table similar to that shown in Figure 3.10, calculate 74
divided by 21 using the hardware described in Figure 3.8. You should show the
contents of each register on each step. Assume both inputs are unsigned 8-bit
integers.

\problembreak

%-----------------------------------------------------------------------------
\subsection{[30] Division Table (Fig.~3.11)}
\hypertarget{prob-3.19}{}
\noindent$\langle$\S3.4$\rangle$\medskip

\noindent Using a table similar to that shown in Figure 3.10, calculate 74
divided by 21 using the hardware described in Figure 3.11. You should show
the contents of each register on each step. This algorithm requires a
slightly different approach than that shown in Figure 3.9. You will want to
think hard about this, do an experiment or two, or else go to the web to
figure out how to make this work correctly. (Hint: one possible solution
involves using the fact that Figure 3.11 implies the remainder register can
be shifted either direction.)

\problembreak

%-----------------------------------------------------------------------------
\subsection{[5] Bit Pattern as Integer}
\hypertarget{prob-3.20}{}
\noindent$\langle$\S3.5$\rangle$\medskip

\noindent What decimal number does the bit pattern \texttt{0x0C000000}
represent if it is a two's complement integer? An unsigned integer?

\problembreak

%-----------------------------------------------------------------------------
\subsection{[10] Bit Pattern as MIPS Instruction}
\hypertarget{prob-3.21}{}
\noindent$\langle$\S3.5$\rangle$\medskip

\noindent If the bit pattern \texttt{0x0C000000} is placed into the
Instruction Register, what MIPS instruction will be executed?

\problembreak

%-----------------------------------------------------------------------------
\subsection{[10] Bit Pattern as Floating Point}
\hypertarget{prob-3.22}{}
\noindent$\langle$\S3.5$\rangle$\medskip

\noindent What decimal number does the bit pattern \texttt{0x0C000000}
represent if it is a floating point number? Use the IEEE 754 standard.

\problembreak

%-----------------------------------------------------------------------------
\subsection{[10] IEEE 754 Single Precision Encoding}
\hypertarget{prob-3.23}{}
\noindent$\langle$\S3.5$\rangle$\medskip

\noindent Write down the binary representation of the decimal number 63.25
assuming the IEEE 754 single precision format.

\problembreak

%-----------------------------------------------------------------------------
\subsection{[10] IEEE 754 Double Precision Encoding}
\hypertarget{prob-3.24}{}
\noindent$\langle$\S3.5$\rangle$\medskip

\noindent Write down the binary representation of the decimal number 63.25
assuming the IEEE 754 double precision format.

\problembreak

%-----------------------------------------------------------------------------
\subsection{[10] IEEE 754 Negative Value Encoding}
\hypertarget{prob-3.25}{}
\noindent$\langle$\S3.5$\rangle$\medskip

\noindent Write down the bit pattern to represent $-1.5625 \times 10^{-1}$
assuming the IEEE 754 single precision format.

\problembreak

%-----------------------------------------------------------------------------
\subsection{[20] DEC PDP-8 Floating Point Format}
\hypertarget{prob-3.26}{}
\noindent$\langle$\S3.5$\rangle$\medskip

\noindent Write down the binary bit pattern to represent
$-1.5625 \times 10^{-1}$ assuming a format similar to that employed by the
DEC PDP-8 (the leftmost 12 bits are the exponent stored as a two's complement
number, and the rightmost 24 bits are the fraction stored as a two's
complement number). No hidden 1 is used. Comment on how the range and
accuracy of this 36-bit pattern compares to the single and double precision
IEEE 754 standards.

\problembreak

%-----------------------------------------------------------------------------
\subsection{[20] IEEE 754-2008 Half Precision Format}
\hypertarget{prob-3.27}{}
\noindent$\langle$\S3.5$\rangle$\medskip

\noindent IEEE 754-2008 contains a half precision that is only 16 bits wide.
The leftmost bit is still the sign bit, the exponent is 5 bits wide and has a
bias of 15, and the mantissa is 10 bits long. A hidden 1 is assumed. Write
down the bit pattern to represent $-1.5625 \times 10^{-1}$ assuming a version
of this format, that uses an excess-16 format to store the exponent. Comment
on how the range and accuracy of this 16-bit floating point format compares
to the single precision IEEE 754 standard.

\problembreak

%-----------------------------------------------------------------------------
\subsection{[20] HP 2114 Floating Point Format}
\hypertarget{prob-3.28}{}
\noindent$\langle$\S3.5$\rangle$\medskip

\noindent The Hewlett-Packard 2114, 2115, and 2116 used a format with the
leftmost 16 bits being the fraction stored in two's complement format,
followed by another 16-bit field which had the leftmost 8 bits as an
extension of the fraction (making the fraction 24 bits long), and the
rightmost 8 bits representing the exponent. However, in an interesting twist,
the exponent was stored in sign-magnitude format with the sign bit on the far
right! Write down the bit pattern to represent $-1.5625 \times 10^{-1}$
assuming this format. No hidden 1 is used. Comment on how the range and
accuracy of this 32-bit pattern compares to the single precision IEEE 754
standard.

\problembreak

%-----------------------------------------------------------------------------
\subsection{[20] Half Precision Addition}
\hypertarget{prob-3.29}{}
\noindent$\langle$\S3.5$\rangle$\medskip

\noindent Calculate the sum of $2.6125 \times 10^{1}$ and
$4.150390625 \times 10^{-1}$ by hand, assuming A and B are stored in the
16-bit half precision described in Exercise 3.27. Assume 1 guard, 1 round
bit, and 1 sticky bit, and round to the nearest even. Show all the steps.

\problembreak

%-----------------------------------------------------------------------------
\subsection{[30] Half Precision Multiplication}
\hypertarget{prob-3.30}{}
\noindent$\langle$\S3.5$\rangle$\medskip

\noindent Calculate the product of $-8.0546875 \times 10^{0}$ and
$-1.79931640625 \times 10^{-1}$ by hand, assuming A and B are stored in the
16-bit half precision format described in Exercise 3.27. Assume 1 guard, 1
round bit, and 1 sticky bit, and round to the nearest even. Show all the
steps; however, as is done in the example in the text, you can do the
multiplication in human-readable format instead of using the techniques
described in Exercises 3.12 through 3.14. Indicate if there is overflow or
underflow. Write your answer in both the 16-bit floating point format
described in Exercise 3.27 and also as a decimal number. How accurate is your
result? How does it compare to the number you get if you do the
multiplication on a calculator?

\problembreak

%-----------------------------------------------------------------------------
\subsection{[30] Half Precision Division}
\hypertarget{prob-3.31}{}
\noindent$\langle$\S3.5$\rangle$\medskip

\noindent Calculate by hand $8.625 \times 10^{1}$ divided by
$-4.875 \times 10^{0}$. Show all the steps necessary to achieve your answer.
Assume there is a guard, a round bit, and a sticky bit, and use them if
necessary. Write the final answer in both the 16-bit floating point format
described in Exercise 3.27 and in decimal and compare the decimal result to
that which you get if you use a calculator.

\problembreak

%-----------------------------------------------------------------------------
\subsection{[20] FP Addition Associativity (Left)}
\hypertarget{prob-3.32}{}
\noindent$\langle$\S3.9$\rangle$\medskip

\noindent Calculate $(3.984375 \times 10^{-1} + 3.4375 \times 10^{-1})
+ 1.771 \times 10^{3}$ by hand, assuming each of the values are stored in the
16-bit half precision format described in Exercise 3.27 (and also described in
the text). Assume 1 guard, 1 round bit, and 1 sticky bit, and round to the
nearest even. Show all the steps, and write your answer in both the 16-bit
floating point format and in decimal.

\problembreak

%-----------------------------------------------------------------------------
\subsection{[20] FP Addition Associativity (Right)}
\hypertarget{prob-3.33}{}
\noindent$\langle$\S3.9$\rangle$\medskip

\noindent Calculate $3.984375 \times 10^{-1} + (3.4375 \times 10^{-1}
+ 1.771 \times 10^{3})$ by hand, assuming each of the values are stored in
the 16-bit half precision format described in Exercise 3.27 (and also
described in the text). Assume 1 guard, 1 round bit, and 1 sticky bit, and
round to the nearest even. Show all the steps and write your answer in both
the 16-bit floating point format and in decimal.

\problembreak

%-----------------------------------------------------------------------------
\subsection{[10] FP Addition Associativity Check}
\hypertarget{prob-3.34}{}
\noindent$\langle$\S3.9$\rangle$\medskip

\noindent Based on your answers to 3.32 and 3.33, does
$(3.984375 \times 10^{-1} + 3.4375 \times 10^{-1}) + 1.771 \times 10^{3}
= 3.984375 \times 10^{-1} + (3.4375 \times 10^{-1} + 1.771 \times 10^{3})$?

\problembreak

%-----------------------------------------------------------------------------
\subsection{[30] FP Multiply Associativity (Left)}
\hypertarget{prob-3.35}{}
\noindent$\langle$\S3.9$\rangle$\medskip

\noindent Calculate $(3.41796875 \times 10^{-3} \times 6.34765625
\times 10^{-3}) \times 1.05625 \times 10^{2}$ by hand, assuming each of the
values are stored in the 16-bit half precision format described in
Exercise 3.27 (and also described in the text). Assume 1 guard, 1 round bit,
and 1 sticky bit, and round to the nearest even. Show all the steps, and
write your answer in both the 16-bit floating point format and in decimal.

\problembreak

%-----------------------------------------------------------------------------
\subsection{[30] FP Multiply Associativity (Right)}
\hypertarget{prob-3.36}{}
\noindent$\langle$\S3.9$\rangle$\medskip

\noindent Calculate $3.41796875 \times 10^{-3} \times (6.34765625
\times 10^{-3} \times 1.05625 \times 10^{2})$ by hand, assuming each of the
values are stored in the 16-bit half precision format described in
Exercise 3.27 (and also described in the text). Assume 1 guard, 1 round bit,
and 1 sticky bit, and round to the nearest even. Show all the steps, and
write your answer in both the 16-bit floating point format and in decimal.

\problembreak

%-----------------------------------------------------------------------------
\subsection{[10] FP Multiply Associativity Check}
\hypertarget{prob-3.37}{}
\noindent$\langle$\S3.9$\rangle$\medskip

\noindent Based on your answers to 3.35 and 3.36, does
$(3.41796875 \times 10^{-3} \times 6.34765625 \times 10^{-3})
\times 1.05625 \times 10^{2} = 3.41796875 \times 10^{-3}
\times (6.34765625 \times 10^{-3} \times 1.05625 \times 10^{2})$?

\problembreak

%-----------------------------------------------------------------------------
\subsection{[30] FP Distributivity (Factored Form)}
\hypertarget{prob-3.38}{}
\noindent$\langle$\S3.9$\rangle$\medskip

\noindent Calculate $1.666015625 \times 10^{0} \times (1.9760 \times 10^{4}
+ (-1.9744 \times 10^{4}))$ by hand, assuming each of the values are stored
in the 16-bit half precision format described in Exercise 3.27 (and also
described in the text). Assume 1 guard, 1 round bit, and 1 sticky bit, and
round to the nearest even. Show all the steps, and write your answer in both
the 16-bit floating point format and in decimal.

\problembreak

%-----------------------------------------------------------------------------
\subsection{[30] FP Distributivity (Expanded Form)}
\hypertarget{prob-3.39}{}
\noindent$\langle$\S3.9$\rangle$\medskip

\noindent Calculate $(1.666015625 \times 10^{0} \times 1.9760 \times 10^{4})
+ (1.666015625 \times 10^{0} \times (-1.9744 \times 10^{4}))$ by hand,
assuming each of the values are stored in the 16-bit half precision format
described in Exercise 3.27 (and also described in the text). Assume 1 guard,
1 round bit, and 1 sticky bit, and round to the nearest even. Show all the
steps, and write your answer in both the 16-bit floating point format and in
decimal.

\problembreak

%-----------------------------------------------------------------------------
\subsection{[10] FP Distributivity Check}
\hypertarget{prob-3.40}{}
\noindent$\langle$\S3.9$\rangle$\medskip

\noindent Based on your answers to 3.38 and 3.39, does
$(1.666015625 \times 10^{0} \times 1.9760 \times 10^{4})
+ (1.666015625 \times 10^{0} \times (-1.9744 \times 10^{4}))
= 1.666015625 \times 10^{0} \times (1.9760 \times 10^{4}
+ (-1.9744 \times 10^{4}))$?

\problembreak

%-----------------------------------------------------------------------------
\subsection{[10] IEEE 754 Representation of $-1/4$}
\hypertarget{prob-3.41}{}
\noindent$\langle$\S3.5$\rangle$\medskip

\noindent Using the IEEE 754 floating point format, write down the bit
pattern that would represent $-1/4$. Can you represent $-1/4$ exactly?

\problembreak

%-----------------------------------------------------------------------------
\subsection{[10] Repeated Addition vs Multiplication}
\hypertarget{prob-3.42}{}
\noindent$\langle$\S3.5$\rangle$\medskip

\noindent What do you get if you add $-1/4$ to itself 4 times? What is
$-1/4 \times 4$? Are they the same? What should they be?

\problembreak

%-----------------------------------------------------------------------------
\subsection{[10] Binary Fraction of $1/3$}
\hypertarget{prob-3.43}{}
\noindent$\langle$\S3.5$\rangle$\medskip

\noindent Write down the bit pattern in the fraction of value $1/3$ assuming a
floating point format that uses binary numbers in the fraction. Assume there
are 24 bits, and you do not need to normalize. Is this representation exact?

\problembreak

%-----------------------------------------------------------------------------
\subsection{[10] BCD Fraction of $1/3$}
\hypertarget{prob-3.44}{}
\noindent$\langle$\S3.5$\rangle$\medskip

\noindent Write down the bit pattern in the fraction of value $1/3$ assuming a
floating point format that uses Binary Coded Decimal (BCD; see Appendix)
(base 10) numbers in the fraction instead of base 2. Assume there are 24
bits, and you do not need to normalize. Is this representation exact?

\problembreak

%-----------------------------------------------------------------------------
\subsection{[10] Base-16 Fraction of $1/3$}
\hypertarget{prob-3.45}{}
\noindent$\langle$\S3.5$\rangle$\medskip

\noindent Write down the bit pattern assuming that we are using 16-bit numbers
in the fraction of value $1/3$ instead of base 2. (Base 16 numbers use the
symbols 0--9 and A--F. Base 15 numbers would use 0--9 and A--E.) Assume there
are 24 bits, and you do not need to normalize. Is this representation exact?

\problembreak

%-----------------------------------------------------------------------------
\subsection{[20] Base-30 Fraction of $1/3$}
\hypertarget{prob-3.46}{}
\noindent$\langle$\S3.5$\rangle$\medskip

\noindent Write down the bit pattern assuming that we are using base 30
numbers in the fraction of value $1/3$ instead of base 2. (Base 16 numbers
use the symbols 0--9 and A--F. Base 30 numbers would use 0--9 and A--T.)
Assume there are 20 bits, and you do not need to normalize. Is this
representation exact?

\problembreak

%-----------------------------------------------------------------------------
\subsection{[45] SIMD FIR Filter Implementation}
\hypertarget{prob-3.47}{}
\noindent$\langle$\S\S3.6, 3.7$\rangle$\medskip

\noindent The following C code implements a four-tap FIR filter on input array
\texttt{sig\_in}. Assume that all arrays are 16-bit fixed-point values.

\begin{verbatim}
    int f[4] = {-1, 2, -1, 2};
    for (i = 3; i < 128; i++)
        sig_out[i] = sig_in[i-3] * f[0] + sig_in[i-2] * f[1]
                   + sig_in[i-1] * f[2] + sig_in[i] * f[3];
\end{verbatim}

\noindent Assume you are to write an optimized implementation of this code in
assembly language on a processor that has SIMD instructions and 128-bit
registers. Without knowing the details of the instruction set, briefly
describe how you would implement this code, maximizing the use of sub-word
operations and minimizing the amount of data that is transferred between
registers and memory. State all your assumptions about the instructions you
use.
